\documentclass[a4paper]{article}
\usepackage{amsfonts}
\usepackage{amsmath}
\usepackage{amssymb}
\usepackage{graphicx}     

\newcommand{\Bgtan}{\operatorname{Bgtan}}
\newcommand{\Bgsin}{\operatorname{Bgsin}}

\begin{document}
  
\noindent \large Auteur: Zaky Souai \\
\noindent \large Studierichting:  Informatica\\
\noindent \large Oefeningenreeks 6A nr. 7.5\\

\bigskip

\normalsize

$ \left\{
\begin{array}{l}
    x_{1} \cdot x_{3} = 36 \\
    x_{2} + x_{4} = 60
\end{array}
\right.$ \\

$ \left\{
\begin{array}{l}
    x_{1} \cdot \left( x_{1} \cdot q^2\right) = 36 \\
    x_{1} \cdot q + x_{1} \cdot q^3 = 60
\end{array}
\right.$ \\

$ \left\{
\begin{array}{l}
    x_{1}^2 \cdot q^2 = 36 \\
    x_{1} \cdot q\left(1 + q^2\right) = 60
\end{array}
\right.$ \\

$ \left\{
\begin{array}{l}
    x_{1} = \pm \dfrac{6}{q} \\
    \left(\pm\dfrac{6}{q} \cdot q\right)(1 + q^2) = 60
\end{array}
\right.$ \\

$ \left\{
\begin{array}{l}
    x_{1} = \pm \dfrac{6}{q} \\
    \pm6\left(1 + q^2\right) = 60
\end{array}
\right.$ \\

$ \left\{
\begin{array}{l}
    q = \pm 3 \\
    x_{1} = \pm 2
\end{array}
\right.$ \\

$\Rightarrow x_{2} = x_{1} \cdot q = 6$\\

$\Rightarrow x_{3} = x_{2} \cdot q = \pm 18$\\

Dus de eerste 3 termen zijn of:\\

$x_{1} = 2$, $x_{2} = 6$, $x_{3} = 18$ \\

of:\\

$x_{1} = -2$, $x_{2} = 6$, $x_{3} = -18$ \\




\begin{thebibliography}{999}
%\bibitem{a} Referentie
\end{thebibliography}
\end{document}